\chapter{Wstęp}
\label{chapt:Wstep}

Dynamiczny rozwój narzędzi do komunikacji zdalnej, zapoczątkowany w latach 90. XX wieku
wraz z upowszechnieniem się Internetu, doprowadził do powstania szerokiej gamy komunikatorów
tekstowych i głosowych. Współczesne rozwiązania, takie jak Microsoft Teams, Discord czy Slack,
oferują rozbudowane funkcjonalności, jednak ich architektura opiera się na centralnych serwerach
zewnętrznych dostawców - oznacza to, że przesyłane dane opuszczają sieć lokalną i mogą być
przetwarzane przez podmioty trzecie.

W środowiskach wymagających zachowania poufności komunikacji, takich jak sieci korporacyjne,
laboratoria badawcze czy instalacje przemysłowe, powyższe rozwiązania mogą stanowić istotne
zagrożenie dla bezpieczeństwa informacji. Aplikacja VT-LAN powstała jako odpowiedź na tę
potrzebę - jest to komunikator tekstowo-głosowy działający wyłącznie w obrębie sieci lokalnej
(LAN), w którym pakiety danych nie są nigdy przekazywane poza jej granice. Kwestia
bezpieczeństwa znalazła odzwierciedlenie również na poziomie doboru bibliotek - jako
warstwę sieciową zastosowano GameNetworkingSockets (GNS), bibliotekę oferującą wbudowane
szyfrowanie transmisji oparte na protokole AES-256-GCM oraz uwierzytelnianie połączeń.
Takie podejście zapewnia ochronę przesyłanych danych nie tylko poprzez ograniczenie ruchu
do sieci lokalnej, lecz także przez szyfrowanie komunikacji na poziomie aplikacji.

Projekt kładzie nacisk na minimalizm i wydajność - aplikacja charakteryzuje się niskim narzutem
na zasoby sprzętowe (CPU, GPU, RAM), prostą obsługą oraz skupieniem na realizacji kluczowych
funkcji komunikacyjnych bez zbędnych dodatków. Kod źródłowy aplikacji udostępniono na licencji
open source, co umożliwia społeczności jego weryfikację, audyt bezpieczeństwa oraz dalszy rozwój.

      
\section{Cel i struktura raportu}
                                       
Celem raportu jest przedstawienie procesu projektowania i implementacji aplikacji
VT-LAN. W raporcie omówiono poszczególne moduły składające się na całość systemu, ze
szczególnym uwzględnieniem zastosowanych narzędzi, bibliotek oraz przyjętych decyzji
projektowych wraz z ich uzasadnieniem.



%\subsection*{Struktura pracy}

%Tutaj opisz krótko, po kolei, wszystkie rozdziały swojej pracy. Oto przykład:

%\begin{quote}
%Rozdział pierwszy zaczynamy (patrz \ref{chapt1:sec:intro}) od opisu wymagań stawianych algorytmom strumieniowym (główne wymaganie to wymaganie użycia %logarytmicznej pamięci względem rozmiaru strumienia), opisujemy pojęcie przesuwnego okna strumienia danych, omawiamy algorytm Bravermana,  ... .\par

%W rozdziale drugim omawiamy ... \par
%...
%\end{quote}

\subsection*{Zastosowane technologie}
Podczas pracy nad aplikacją VT-LAN zespół korzystał z następujących technologii:
\begin{enumerate}
	\item \textbf{Język C++} - główny i jedyny język wykorzystywany do implementacji
	logiki aplikacji. Odpowiada zarówno za programowanie interfejsu użytkownika,
	jak i za obsługę komunikacji sieciowej, przechwytywanie oraz odtwarzanie dźwięku.
	Wybór C++ podyktowany był wymaganiami wydajnościowymi projektu - język ten
	zapewnia niski narzut pamięciowy i pełną kontrolę nad zasobami systemowymi,
	co jest kluczowe przy przetwarzaniu strumieni audio w czasie rzeczywistym.
	
	\item \textbf{Microsoft Visual Studio 2022} - zintegrowane środowisko deweloperskie
	(IDE) wraz z kompilatorem MSVC, wykorzystywane jako główne narzędzie do budowania
	i debugowania projektu. Dostarcza zaawansowane narzędzia diagnostyczne \cite{msvc_profiler},
	w tym debugger z obsługą wielowątkowości oraz zintegrowany pakiet \textit{Diagnostic Tools}
	umożliwiający profilowanie zużycia CPU, wykonywanie migawek sterty (ang.~\textit{heap
	snapshots}) i lokalizowanie wycieków pamięci w czasie rzeczywistym -- bez konieczności
	modyfikacji kodu produkcyjnego.
	
	\item \textbf{Premake5} - system budowania projektów oparty na skryptach Lua.
	Umożliwia generowanie plików rozwiązania (\texttt{.sln}) dla Visual Studio
	z pojedynczego, czytelnego pliku konfiguracyjnego. Zastosowanie Premake5
	zamiast bezpośredniej edycji plików projektu upraszcza zarządzanie zależnościami
	oraz ułatwia konfigurację środowiska członkom zespołu - wszelkie zmiany
	struktury projektu zawarte są w plikach Lua i automatycznie odzwierciedlane
	przy kolejnym generowaniu rozwiązania.
	
	\item \textbf{GitHub} - platforma do współdzielenia kodu źródłowego oparta na systemie
	kontroli wersji Git. Służyła jako centralne narzędzie do koordynacji prac zespołu
	- przeglądania zmian w kodzie, udostępniania aktualizacji oraz śledzenia zadań
	i błędów za pomocą systemu zgłoszeń. Repozytorium pełni również rolę miejsca
	dystrybucji projektu - użytkownicy mają możliwość wglądu w kod źródłowy aplikacji,
	a także pobrania skompilowanej wersji wykonywalnej.
	
	\item \textbf{Wireshark} \cite{wireshark_doc} - analizator protokołów sieciowych
	(ang.~\textit{packet sniffer}) wykorzystywany podczas prac diagnostycznych i kontroli
	jakości (QA). Umożliwiał przechwytywanie i inspekcję ruchu sieciowego generowanego
	przez aplikację, co pozwoliło zweryfikować poprawność działania protokołów UDP i TCP/GNS,
	skuteczność szyfrowania pakietów audio oraz zachowanie mechanizmu detekcji aktywności
	głosowej (VAD) -- szczegóły opisano w rozdziale~\ref{ch:profilowanie_debugging}.

	\item \textbf{Biblioteki firm trzecich} - projekt korzysta z następujących bibliotek:
	\begin{itemize}
		\item \textbf{SDL3} \cite{SDL3} - wieloplatformowa biblioteka obsługi okna aplikacji,
		renderowania interfejsu użytkownika oraz systemu audio. W zakresie dźwięku
		SDL3 dostarcza niskopoziomowe API do otwierania urządzeń wejściowych
		i wyjściowych, przechwytywania próbek audio z mikrofonu oraz ich odtwarzania
		przez głośniki. Do miksowania i nakładania strumieni audio aplikacja korzysta
		z biblioteki \textbf{SDL\_mixer} \cite{SDLMixer} - oficjalnego rozszerzenia SDL3
		umożliwiającego jednoczesne odtwarzanie wielu kanałów audio z regulacją głośności.
		
		\item \textbf{Dear ImGui} \cite{ImGui} - biblioteka interfejsu użytkownika
		dla języka C++. Odpowiada za renderowanie
		wszystkich elementów graficznych aplikacji - okien, przycisków, pól tekstowych,
		list oraz paneli czatu i głosowego. Charakteryzuje się minimalną liczbą
		zależności i niskim narzutem wydajnościowym, co czyni ją szczególnie
		odpowiednią dla aplikacji czasu rzeczywistego. Integruje się bezpośrednio
		z SDL3 jako backendem okienkowym i wejściowym.
		
		\item \textbf{Opus} \cite{OpusRFC} - otwarty kodek audio o niskim opóźnieniu,
		zoptymalizowany pod transmisję mowy, ze zmiennym bitrate dostosowującym się
		do warunków sieci. Bezpośrednio integruje się z buforami próbek pobieranymi
		przez SDL3, kompresując dane przed wysyłką przez GNS i dekompresując je
		po stronie odbiorcy. Zastosowanie Opus pozwala znacząco zredukować
		przepustowość wymaganą do transmisji głosu przy zachowaniu wysokiej jakości
		dźwięku.
		
		\item \textbf{GameNetworkingSockets (GNS)} \cite{GNS} - biblioteka sieciowa
		firmy Valve zapewniająca niezawodną, szyfrowaną komunikację opartą na protokole
		UDP. Oferuje mechanizmy detekcji opóźnień, kontrolę przepływu oraz opcjonalne
		gwarancje dostarczenia pakietów, co czyni ją dobrze dopasowaną zarówno
		do komunikacji klient-serwer, jak i peer-to-peer w sieci LAN. Szyfrowanie
		realizowane jest przy użyciu algorytmu AES-GCM-256 per pakiet,
		z wymianą kluczy opartą na Curve25519 \cite{GNS}. GNS korzysta wewnętrznie
		z następujących bibliotek:
		\begin{itemize}
			\item \textbf{OpenSSL} \cite{OpenSSL} - biblioteka kryptograficzna
			zapewniająca szyfrowanie transmisji oraz uwierzytelnianie połączeń,
			\item \textbf{protobuf} \cite{protobuf} - biblioteka Google Protocol Buffers
			wykorzystywana do serializacji wewnętrznych komunikatów sterujących
			protokołu GNS.
		\end{itemize}
	\end{itemize}
\end{enumerate}

\section{Wykorzystanie AI}

W trakcie realizacji niniejszej pracy projektowej zespół korzystał z narzędzi
generatywnej sztucznej inteligencji zgodnie z wytycznymi Wydziału Informatyki
i Telekomunikacji Politechniki Wrocławskiej \cite{WITWytyczneSI}.
Wykorzystywane były następujące narzędzia:

\begin{itemize}
	\item \textbf{Claude (Anthropic, wersja Pro)},
	\item \textbf{ChatGPT (OpenAI, wersja Pro)},
	\item \textbf{Gemini (Google)}.
\end{itemize}

Narzędzia te były stosowane wyłącznie jako pomoc techniczna i redakcyjna,
w następujących obszarach:

\begin{itemize}
	\item \textbf{Wspomaganie implementacji} -- generowanie fragmentów kodu,
	podpowiedzi dotyczących struktury rozwiązań oraz debugowanie błędów.
	Fragmenty kodu powstałe przy udziale narzędzi AI zostały oznaczone
	stosownymi komentarzami bezpośrednio w plikach źródłowych projektu,
	
	\item \textbf{Wyszukiwanie informacji o bibliotekach i technologiach} --
	zbieranie informacji na temat dostępnych bibliotek, przyjętych standardów
	oraz rozwiązań stosowanych w branży (m.in. GameNetworkingSockets, Opus, SDL3),
	
	\item \textbf{Korekta stylistyczna i językowa tekstu} -- poprawa poprawności
	gramatycznej, spójności stylu oraz dostosowanie języka do wymogów raportu.
	
	\item \textbf{Wsparcie przy diagnostyce i profilowaniu} --
	w rozdziale dotyczącym profilowania i debuggingu wykorzystano model
	Gemini (Google) jako asystenta przy ustrukturyzowaniu dokumentu,
	konsultacjach metodyki testowej dotyczącej filtrów diagnostycznych
	Wireshark dla strumieni audio opartych na kodeku Opus, oraz przy
	formułowaniu składni bloków kodu C++ związanych z mechanizmami
	synchronizacji wątków. Wszelkie wyniki liczbowe, wykresy wydajnościowe
	oraz analiza błędów opierają się na rzeczywistych testach
	przeprowadzonych w środowisku laboratoryjnym.
\end{itemize}

Autorzy pracy zachowali pełną odpowiedzialność za merytoryczną poprawność
wszystkich zawartych w niej treści. Każda informacja pozyskana przy udziale
narzędzi AI była samodzielnie weryfikowana, w szczególności w zakresie
poprawności bibliograficznej i technicznej. Narzędzia AI nie były wykorzystywane
do generowania całych rozdziałów ani struktury pracy.

\section{Opłacalność projektu}

VT-LAN jest projektem open source udostępnianym nieodpłatnie, którego model
finansowy opiera się na zerowym koszcie utrzymania po stronie autorów.
Całość infrastruktury działa wyłącznie na maszynach użytkowników w sieci lokalnej
- nie istnieje żaden centralny serwer, który wymagałby hostowania, opłacania
ani administrowania. Koszt utrzymania projektu po stronie twórców jest zatem
zerowy.

Trwałość projektu zapewniana jest przez społeczność - model open source
umożliwia każdemu użytkownikowi zgłaszanie poprawek, łatanie usterek oraz
rozwijanie nowych funkcjonalności w formie tzw. patchy. Dzięki temu jakość
aplikacji może rosnąć bez ponoszenia przez autorów kosztów deweloperskich.
Dodatkowym, dobrowolnym źródłem finansowania mogą być dotacje przekazywane
przez użytkowników (np. poprzez platformy typu GitHub Sponsors lub Open
Collective), co jest powszechnie stosowaną praktyką w projektach open source.

Biorąc pod uwagę powyższe, projekt cechuje się pełną opłacalnością przy
zerowym progu wejścia - zarówno dla autorów, jak i dla użytkowników końcowych.

\section{Użyteczność społeczna produktu}
\label{sec:użytek}
Rosnąca świadomość zagrożeń związanych z prywatnością danych w komunikatorach
internetowych czyni VT-LAN odpowiedzią na realną potrzebę społeczną. Popularne
platformy komunikacyjne, takie jak Discord, Teams czy Slack, przesyłają dane
głosowe i tekstowe przez zewnętrzne serwery swoich dostawców, co oznacza
potencjalną ekspozycję tych danych na osoby trzecie.

Trafnym przykładem ilustrującym tę obawę jest kontrowersja wokół Discorda
z początku 2026 roku. Platforma ogłosiła wprowadzenie globalnego systemu
weryfikacji wieku, wymagającego od użytkowników przesłania skanu dowodu
tożsamości lub wykonania skanu twarzy. Decyzja ta wywołała falę krytyki
ze strony społeczności i organizacji zajmujących się ochroną prywatności
- tym bardziej, że we wrześniu 2025 roku doszło do wycieku zdjęć dokumentów
tożsamości nawet 70\,000 użytkowników Discorda wskutek naruszenia bezpieczeństwa
u zewnętrznego dostawcy usług \cite{DiscordEFF}. Skala reakcji była na tyle
duża, że Discord oficjalnie opóźnił wdrożenie systemu do drugiej połowy 2026
roku \cite{DiscordPostpone}.

Incydent ten uwidocznił fundamentalny problem centralnych komunikatorów:
dane użytkowników znajdują się poza ich kontrolą. VT-LAN eliminuje ten problem
u podstaw - ruch sieciowy nigdy nie opuszcza sieci lokalnej, a szyfrowanie
na poziomie aplikacji (AES-GCM-256) dodatkowo chroni przesyłane treści.
Aplikacja jest szczególnie użyteczna w środowiskach wymagających poufności,
takich jak sieci korporacyjne, laboratoria, instytucje edukacyjne czy
domowe sieci prywatne.

Jako projekt open source VT-LAN wpisuje się ponadto w ideę transparentności
oprogramowania - każdy użytkownik może zweryfikować kod źródłowy i upewnić
się, że aplikacja działa zgodnie z deklarowanymi zasadami. Repozytorium
dostępne publicznie w serwisie GitHub umożliwia pobranie dowolnej wersji
projektu wraz z pełnym kodem źródłowym i samodzielne skompilowanie go na
własnej maszynie. Dzięki temu użytkownicy nie są uzależnieni od decyzji
autorów - w przeciwieństwie do zamkniętych platform, takich jak Discord,
żadna przyszła aktualizacja nie może bez wiedzy społeczności wprowadzić
mechanizmów weryfikacji tożsamości ani przetwarzania danych biometrycznych.

\section{Plan finansowy projektu}

Ze względu na otwartoźródłowy charakter aplikacji oraz jej architekturę
opartą wyłącznie na infrastrukturze użytkownika, VT-LAN nie generuje
żadnych stałych kosztów operacyjnych po stronie autorów. Poniżej
przedstawiono zestawienie kosztów i potencjalnych przychodów projektu.

\subsection*{Koszty}

\begin{itemize}
	\item \textbf{Hosting i infrastruktura serwerowa} - brak. Aplikacja
	działa wyłącznie na maszynach użytkowników w sieci LAN; autorzy nie
	ponoszą żadnych kosztów serwerowych.
	\item \textbf{Licencje} - brak. Wszystkie zastosowane biblioteki
	(SDL3, GameNetworkingSockets, Opus, protobuf) udostępniane są na
	licencjach open source, niewymagających żadnych opłat.
	\item \textbf{Dystrybucja} - brak. Projekt dystrybuowany jest
	bezpłatnie za pośrednictwem repozytorium GitHub.
	\item \textbf{Utrzymanie} - pokrywane społecznie w modelu open source
	poprzez zgłoszenia błędów, patche oraz pull requesty. Autorzy mogą
	hobbystycznie, w miarę dostępnego czasu, weryfikować i zatwierdzać
	wkład społeczności lub samodzielnie rozwijać projekt.
\end{itemize}

\subsection*{Przychody}

\begin{itemize}
	\item \textbf{Dobrowolne dotacje użytkowników} -- projekt może korzystać
	z platform umożliwiających dobrowolne wspieranie finansowe twórców
	open source, takich jak GitHub Sponsors lub Open Collective.
	Przychody z tego źródła mają charakter nieregularny i zależą
	od zaangażowania społeczności.
\end{itemize}

Projekt nie jest nastawiony na generowanie zysku. Jego wartość mierzona jest
społecznym zasięgiem, liczbą aktywnych użytkowników oraz wkładem społeczności
w jego rozwój.

\section{Przegląd rozwiązań dostępnych na rynku}
\label{sec:przeglad_rynku}

Na rynku komunikatorów głosowych i tekstowych funkcjonuje kilka ugruntowanych
rozwiązań, które można uznać za punkty odniesienia dla projektu VT-LAN.
Poniżej przedstawiono ich charakterystykę wraz z oceną pod kątem prywatności,
architektury sieciowej i otwartości kodu.

\subsection*{Discord}

Discord \cite{DiscordEFF} to jedna z najpopularniejszych platform komunikacyjnych,
oferująca czat tekstowy, głosowy oraz wideo. Aplikacja jest
bezpłatna w wersji podstawowej, z opcjonalną subskrypcją Discord Nitro.

Z punktu widzenia prywatności Discord jest rozwiązaniem centralnym - cały
ruch sieciowy przechodzi przez serwery firmy, co oznacza brak kontroli
użytkownika nad przesyłanymi danymi. Kontrowersje wokół próby wprowadzenia
przez platformę weryfikacji wieku poprzez skanowanie dokumentu tożsamości
i danych biometrycznych, opisane szczegółowo w sekcji~\ref{sec:użytek},
uwidoczniły systemowe ryzyko wynikające z centralizacji danych
\cite{DiscordPostpone}.

\subsection*{TeamSpeak}

TeamSpeak \cite{TeamSpeak} to wieloletnia platforma VoIP skupiona przede
wszystkim na komunikacji głosowej, szeroko stosowana w środowiskach graczy
oraz w zastosowaniach profesjonalnych. Wyróżnia się niskim opóźnieniem
transmisji oraz możliwością uruchomienia własnego serwera, co daje
administratorowi pełną kontrolę nad infrastrukturą. Wersja podstawowa
jest bezpłatna, natomiast utrzymanie własnego serwera wiąże się z kosztami
hostingu. Kod źródłowy pozostaje zamknięty.

\subsection*{Mumble}

Mumble \cite{Mumble} jest otwartoźródłowym komunikatorem głosowym
o niskim opóźnieniu, działającym w modelu klient-serwer. Serwer (Murmur)
może być uruchomiony w sieci lokalnej, co eliminuje konieczność przesyłania
danych przez serwery zewnętrzne. Aplikacja jest w pełni bezpłatna i dostępna na licencji
open source. Słabą stroną Mumble jest brak natywnej obsługi przesyłania plików.

\subsection*{Microsoft Teams}

Microsoft Teams to korporacyjna platforma komunikacyjna zintegrowana
z ekosystemem Microsoft 365. Oferuje czat tekstowy, połączenia głosowe
i wideo, a także współdzielenie plików i zaawansowane narzędzia pracy
zespołowej. Jest rozwiązaniem wyłącznie chmurowym - dane przetwarzane są
na serwerach Microsoft, co w środowiskach o podwyższonych wymaganiach
bezpieczeństwa może stanowić istotne ograniczenie. Produkt jest płatny
w ramach subskrypcji Microsoft 365.

\subsection*{Podsumowanie i pozycja VT-LAN}
VT-LAN zajmuje unikalną niszę wśród powyższych rozwiązań. Jest jedynym
komunikatorem zaprojektowanym z myślą o działaniu wyłącznie w sieci LAN,
łączącym czat tekstowy, transmisję głosową i przesyłanie plików w jednej
lekkiej aplikacji open source, niewymagającej żadnej infrastruktury
zewnętrznej ani kont użytkowników.

\begin{table}[H]
	\centering
	\caption{Porównanie komunikatorów dostępnych na rynku}
	\begin{tabularx}{\textwidth}{|l|X|X|X|X|X|}
		\hline
		\textbf{Cecha} & \textbf{Discord} & \textbf{TeamSpeak} & \textbf{Mumble}
		& \textbf{MS Teams} & \textbf{VT-LAN} \\
		\hline
		Czat tekstowy       & Tak  & Ograniczony & Ograniczony & Tak  & Tak \\
		\hline
		Głos (VoIP)         & Tak  & Tak         & Tak         & Tak  & Tak \\
		\hline
		Przesyłanie plików  & Tak  & Nie         & Nie         & Tak  & Tak \\
		\hline
		Tylko LAN           & Nie  & Nie         & Możliwe     & Nie  & Tak \\
		\hline
		Open source         & Nie  & Nie         & Tak         & Nie  & Tak \\
		\hline
		Własny serwer       & Nie  & Tak         & Tak         & Nie  & Tak \\
		\hline
		Bezpłatny           & Częściowo & Częściowo & Tak    & Nie  & Tak \\
		\hline
		Narzut zasobów      & Wysoki & Niski     & Niski       & Wysoki & Niski \\
		\hline
	\end{tabularx}
\end{table}
